% Generated from docs/book.pdf. Review equations and figures against the source PDF.
\chapter{Kernel Machines}
\label{chap:13-kernel-machines}

Kernel machines are maximum margin methods that allow the model to be written as a sum of the influences of a subset of the training instances. These influences are given by application-specific similarity kernels, and we discuss “kernelized” classification, regression, outlier detection, and dimensionality reduction, as well as how to choose and use kernels.

\section{Introduction}

We now discuss a different approach for linear classification and regression. We should not be surprised to have so many different methods even for the simple case of a linear model. Each learning algorithm has a different inductive bias, makes different assumptions, and defines a different objective function and thus may find a different linear model. The model that we will discuss in this chapter, called the support vector machine (SVM), and later generalized under the name kernel machine, has been popular in recent years for a number of reasons:

1. It is a discriminant-based method and uses Vapnik’s principle to never solve a more complex problem as a first step before the actual problem (Vapnik 1995). For example, in classification, when the task is to learn the discriminant, it is not necessary to estimate where the class densities p(x|Ci) or the exact posterior probability values P(Ci|x); we only need to estimate where the class boundaries lie, that is, x where P(Ci|x) = P(Cj|x). Similarly, for outlier detection, we do not need to estimate the full density p(x); we only need to find the boundary separating those x that have low p(x), that is, x where p(x) < $\theta$, for some threshold $\theta$ $\in$(0, 1).

2. After training, the parameter of the linear model, the weight vector, can be written down in terms of a subset of the training set, which are the so-called support vectors. In classification, these are the cases that are close to the boundary and as such, knowing them allows knowledge extraction: those are the uncertain or erroneous cases that lie in the vicinity of the boundary between two classes. Their number gives us an estimate of the generalization error, and, as we see below, being able to write the model parameter in terms of a set of instances allows kernelization.

3. As we will see shortly, the output is written as a sum of the influences of support vectors and these are given by kernel functions that are application-specific measures of similarity between data instances. Previously, we talked about nonlinear basis functions allowing us to map the input to another space where a linear (smooth) solution is possible; the kernel function uses the same idea.

4. Typically in most learning algorithms, data points are represented as vectors, and either dot product (as in the multilayer perceptrons) or Euclidean distance (as in radial basis function networks) is used. A kernel function allows us to go beyond that. For example, G1 and G2 may be two graphs and K(G1, G2) may correspond to the number of shared paths, which we can calculate without needing to represent G1 or G2 explicitly as vectors.

5. Kernel-based algorithms are formulated as convex optimization problems, and there is a single optimum that we can solve for analytically. Therefore we are no longer bothered with heuristics for learning rates, initializations, checking for convergence, and such. Of course, this does not mean that we do not have any hyperparameters for model selection; we do---any method needs them, to match the algorithm to the data at hand.

We start our discussion with the case of classification, and then generalize to regression, outlier (novelty) detection, and then dimensionality reduction. We see that in all cases basically we have the similar quadratic program template to maximize the separability, or margin, of instances subject to a constraint of the smoothness of solution. Solving for it, we get the support vectors. The kernel function defines the space according to its notion of similarity and a kernel function is good if we have better separation in its corresponding space.

\section{Optimal Separating Hyperplane}

Let us start again with two classes and use labels $-$1/ + 1 for the two classes. The sample is X = \{xt, rt\} where rt = +1 if xt $\in$C1 and rt = $-$1 if xt $\in$C2. We would like to find w and w0 such that

wT xt + w0 $\geq$+1 for rt = +1 wT xt + w0 $\leq$$-$1 for rt = $-$1

which can be rewritten as

(13.1) rt(wTxt + w0) $\geq$+1

Note that we do not simply require

rt(wTxt + w0) $\geq$0

Not only do we want the instances to be on the right side of the hyperplane, but we also want them some distance away, for better generalization. The distance from the hyperplane to the instances closest to it margin on either side is called the margin, which we want to maximize for best generalization. Very early on, in section 2.1, we talked about the concept of the margin when we were talking about fitting a rectangle, and we said that it is better to take a rectangle halfway between S and G, to get a breathing space. This is so that in case noise shifts a test instance slightly, it will still be on the right side of the boundary. optimal separating Similarly, now that we are using the hypothesis class of lines, the optihyperplane mal separating hyperplane is the one that maximizes the margin. We remember from section 10.3 that the distance of xt to the discriminant is

|wTxt + w0| $\parallel$w$\parallel$

which, when rt $\in$\{$-$1, +1\}, can be written as

rt(wTxt + w0) $\parallel$w$\parallel$

and we would like this to be at least some value $\rho$:

rt(wTxt + w0) (13.2) $\geq$$\rho$, $\forall$t $\parallel$w$\parallel$

We would like to maximize $\rho$ but there are an infinite number of solutions that we can get by scaling w and for a unique solution, we fix $\rho$$\parallel$w$\parallel$= 1 and thus, to maximize the margin, we minimize $\parallel$w$\parallel$. The task can therefore be defined (see Cortes and Vapnik 1995; Vapnik 1995) as to

(13.3) min 2$\parallel$w$\parallel$2 subject to rt(wT xt + w0) $\geq$+1, $\forall$t

This is a standard quadratic optimization problem, whose complexity depends on d, and it can be solved directly to find w and w0. Then, on both sides of the hyperplane, there will be instances that are 1/$\parallel$w$\parallel$away from the hyperplane and the total margin will be 2/$\parallel$w$\parallel$. We saw in section 10.2 that if the problem is not linearly separable, instead of fitting a nonlinear function, one trick is to map the problem to a new space by using nonlinear basis functions. It is generally the case that this new space has many more dimensions than the original space, and, in such a case, we are interested in a method whose complexity does not depend on the input dimensionality. In finding the optimal hyperplane, we can convert the optimization problem to a form whose complexity depends on N, the number of training instances, and not on d. Another advantage of this new formulation is that it will allow us to rewrite the basis functions in terms of kernel functions, as we will see in section 13.5. To get the new formulation, we first write equation 13.3 as an unconstrained problem using Lagrange multipliers $\alpha$t:

N Lp = 2$\parallel$w$\parallel$2 $-$ $\alpha$t[rt(wTxt + w0) $-$1] t=1 (13.4) = 2$\parallel$w$\parallel$2 $-$ $\alpha$trt(wT xt + w0) + $\alpha$t t t

This should be minimized with respect to w, w0 and maximized with respect to $\alpha$t $\geq$0. The saddle point gives the solution. This is a convex quadratic optimization problem because the main term is convex and the linear constraints are also convex. Therefore, we can equivalently solve the dual problem, making use of the Karush-Kuhn- Tucker conditions. The dual is to maximize Lp with respect to $\alpha$t, subject to the constraints that the gradient of Lp with respect to w and w0 are 0

and also that $\alpha$t $\geq$0:

$\partial$Lp (13.5) = 0 $\Rightarrow$ w = $\alpha$trtxt $\partial$w t $\partial$Lp (13.6) = 0 $\Rightarrow$ $\alpha$trt = 0 $\partial$w0 t

Plugging these into equation 13.4, we get the dual

Ld = 2(wT w) $-$wT $\alpha$trtxt $-$w0 $\alpha$trt + $\alpha$t t t t = $-$1 w) + $\alpha$t 2(wT t (13.7) = $-$1 $\alpha$t$\alpha$srtrs(xt)T xs + $\alpha$t 2 t s t

which we maximize with respect to $\alpha$t only, subject to the constraints

$\alpha$trt = 0, and $\alpha$t $\geq$0, $\forall$t t

This can be solved using quadratic optimization methods. The size of the dual depends on N, sample size, and not on d, the input dimensionality. The upper bound for time complexity is O(N3), and the upper bound for space complexity is O(N2). Once we solve for $\alpha$t, we see that though there are N of them, most vanish with $\alpha$t = 0 and only a small percentage have $\alpha$t > 0. The set of xt whose $\alpha$t > 0 are the support vectors, and as we see in equation 13.5, w is written as the weighted sum of these training instances that are selected as the support vectors. These are the xt that satisfy

rt(wTxt + w0) = 1

and lie on the margin. We can use this fact to calculate w0 from any support vector as

(13.8) w0 = rt $-$wTxt

For numerical stability, it is advised that this be done for all support vectors and an average be taken. The discriminant thus found is called support vector the support vector machine (SVM) (see figure 13.1). machine The majority of the $\alpha$t are 0, for which rt(wTxt + w0) > 1. These are the xt that lie more than sufficiently away from the discriminant,

1.5

0.5

$-$1 0 0.5 1 1.5 2

Figure 13.1 For a two-class problem where the instances of the classes are shown by plus signs and dots, the thick line is the boundary and the dashed lines define the margins on either side. Circled instances are the support vectors.

and they have no effect on the hyperplane. The instances that are not support vectors carry no information; even if any subset of them are removed, we would still get the same solution. From this perspective, the SVM algorithm can be likened to the condensed nearest neighbor algorithm (section 8.5), which stores only the instances neighboring (and hence constraining) the class discriminant. Being a discriminant-based method, the SVM cares only about the instances close to the boundary and discards those that lie in the interior. Using this idea, it is possible to use a simpler classifier before the SVM to filter out a large portion of such instances, thereby decreasing the complexity of the optimization step of the SVM (exercise 1). During testing, we do not enforce a margin. We calculate g(x) = wT x + w0, and choose according to the sign of g(x):

Choose C1 if g(x) > 0 and C2 otherwise

\section{The Nonseparable Case: Soft Margin Hyperplane}

If the data is not linearly separable, the algorithm we discussed earlier will not work. In such a case, if the two classes are not linearly separable such that there is no hyperplane to separate them, we look for the one slack variables that incurs the least error. We define slack variables, $\xi$t $\geq$0, which store the deviation from the margin. There are two types of deviation: An instance may lie on the wrong side of the hyperplane and be misclassified. Or, it may be on the right side but may lie in the margin, namely, not sufficiently away from the hyperplane. Relaxing equation 13.1, we require

(13.9) rt(wTxt + w0) $\geq$1 $-$$\xi$t

If $\xi$t = 0, there is no problem with xt. If 0 < $\xi$t < 1, xt is correctly classified but in the margin. If $\xi$t $\geq$1, xt is misclassified (see figure 13.2). The number of misclassifications is \#\{$\xi$t > 1\}, and the number of nonsoft error separable points is \#\{$\xi$t > 0\}. We define soft error as

$\xi$t t

and add this as a penalty term:

(13.10) Lp = 2$\parallel$w$\parallel$2 + C $\xi$t t

subject to the constraint of equation 13.9. C is the penalty factor as in any regularization scheme trading offcomplexity, as measured by the L2 norm of the weight vector (similar to weight decay in multilayer perceptrons; see section 11.9), and data misfit, as measured by the number of nonseparable points. Note that we are penalizing not only the misclassified points but also the ones in the margin for better generalization, though these latter would be correctly classified during testing. Adding the constraints, the Lagrangian of equation 13.4 then becomes

(13.11) Lp = 2$\parallel$w$\parallel$2 + C $\xi$t $-$ $\alpha$t[rt(wT xt + w0) $-$1 + $\xi$t] $-$ $\mu$t$\xi$t t t t

where $\mu$t are the new Lagrange parameters to guarantee the positivity of $\xi$t. When we take the derivatives with respect to the parameters and set them to 0, we get:

$\partial$Lp (13.12) = w $-$ $\alpha$trtxt = 0 $\Rightarrow$w = $\alpha$trtxt $\partial$w t t

(a)

1.5

(b) 1 (c)

(d) 0.5

$-$1 1

0 0.5 1 1.5 2

Figure 13.2 In classifying an instance, there are four possible cases: In (a), the instance is on the correct side and far away from the margin; r tg(xt) > 1, $\xi$t = 0. In (b), $\xi$t = 0; it is on the right side and on the margin. In (c), $\xi$t = 1 $-$g(xt), 0 < $\xi$ < 1; it is on the right side but is in the margin and not sufficiently away. In (d), $\xi$t = 1 + g(xt) > 1; it is on the wrong side---this is a misclassification. All cases except (a) are support vectors. In terms of the dual variable, in (a), $\alpha$t = 0; in (b), $\alpha$t < C; in (c) and (d), $\alpha$t = C.

$\partial$Lp (13.13) = $\alpha$trt = 0 $\partial$w0 t $\partial$Lp (13.14) = C $-$$\alpha$t $-$$\mu$t = 0 $\partial$$\xi$t

Since $\mu$t $\geq$0, this last implies that 0 $\leq$$\alpha$t $\leq$C. Plugging these into equation 13.11, we get the dual that we maximize with respect to $\alpha$t:

(13.15) Ld = $\alpha$t $-$1 $\alpha$t$\alpha$srtrs(xt)T xs 2 t t s

subject to

$\alpha$trt = 0 and 0 $\leq$$\alpha$t $\leq$C, $\forall$t t

Solving this, we see that as in the separable case, instances that lie on the correct side of the boundary with sufficient margin vanish with their $\alpha$t = 0 (see figure 13.2). The support vectors have their $\alpha$t > 0 and they define w, as given in equation 13.12. Of these, those whose $\alpha$t < C are the ones that are on the margin, and we can use them to calculate w0; they have $\xi$t = 0 and satisfy rt(wTxt + w0) = 1. Again, it is better to take an average over these w0 estimates. Those instances that are in the margin or misclassified have their $\alpha$t = C. The nonseparable instances that we store as support vectors are the instances that we would have trouble correctly classifying if they were not in the training set; they would either be misclassified or classified correctly but not with enough confidence. We can say that the number of support vectors is an upper-bound estimate for the expected number of errors. And, actually, Vapnik (1995) has shown that the expected test error rate is EN[P(error)] $\leq$EN[\# of support vectors] N where EN[$\cdot$] denotes expectation over training sets of size N. The nice implication of this is that it shows that the error rate depends on the number of support vectors and not on the input dimensionality. Equation 13.9 implies that we define error if the instance is on the hinge loss wrong side or if the margin is less than 1. This is called the hinge loss. If yt = wT xt + w0 is the output and rt is the desired output, hinge loss is defined as 0 if ytrt $\geq$1 (13.16) Lhinge(yt, rt) = 1 $-$ytrt otherwise

In figure 13.3, we compare hinge loss with 0/1 loss, squared error, and cross-entropy. We see that different from 0/1 loss, hinge loss also penalizes instances in the margin even though they may be on the correct side, and the loss increases linearly as the instance moves away on the wrong side. This is different from the squared loss that therefore is not as robust as the hinge loss. We see that cross-entropy minimized in logistic discrimination (section 10.7) or by the linear perceptron (section 11.3), is a good continuous approximation to the hinge loss. C of equation 13.10 is the regularization parameter fine-tuned using cross-validation. It defines the trade-offbetween margin maximization and error minimization: If it is too large, we have a high penalty for nonseparable points, and we may store many support vectors and overfit.

= rt 5 squared error for loss 3 hinge loss

2 cross entropy

0/1 loss $-$2 $-$1 0 1 2 yt

Figure 13.3 Comparison of different loss functions for r t = 1: 0/1 loss is 0 if yt = 1, 1 otherwise. Hinge loss is 0 if yt > 1, 1 $-$yt otherwise. Squared error is (1 $-$yt)2. Cross-entropy is log(1/(1 + exp($-$yt))).

If it is too small, we may find too simple solutions that underfit. Typically, one chooses from [10$-$6, 10$-$5, . . . , 10+5, 10+6] in the log scale by looking at the accuracy on a validation set.

\section{nu-SVM}

There is another, equivalent formulation of the soft margin hyperplane that uses a parameter $\nu$ $\in$[0, 1] instead of C (Schölkopf et al. 2000). The objective function is

1 1 (13.17) min 2$\parallel$w$\parallel$2 $-$$\nu$$\rho$ + N $\xi$t t

subject to

(13.18) rt(wTxt + w0) $\geq$$\rho$ $-$$\xi$t, $\xi$t $\geq$0, $\rho$ $\geq$0

$\rho$ is a new parameter that is a variable of the optimization problem and scales the margin: the margin is now 2$\rho$/$\parallel$w$\parallel$. $\nu$ has been shown to be a

lower bound on the fraction of support vectors and an upper bound on the fraction of instances having margin errors ( t \#\{$\xi$t > 0\}). The dual is (13.19) Ld = $-$1 $\alpha$t$\alpha$srtrs(xt)Txs 2 t s

subject to $\alpha$trt = 0, 0 $\leq$$\alpha$t $\leq$1 , $\alpha$t $\leq$$\nu$ N t t

When we compare equation 13.19 with equation 13.15, we see that the term t $\alpha$t no longer appears in the objective function but is now a constraint. By playing with $\nu$, we can control the fraction of support vectors, and this is advocated to be more intuitive than playing with C.

\section{Kernel Trick}

Section 10.2 demonstrated that if the problem is nonlinear, instead of trying to fit a nonlinear model, we can map the problem to a new space by doing a nonlinear transformation using suitably chosen basis functions and then use a linear model in this new space. The linear model in the new space corresponds to a nonlinear model in the original space. This approach can be used in both classification and regression problems, and in the special case of classification, it can be used with any scheme. In the particular case of support vector machines, it leads to certain simplifications that we now discuss. Let us say we have the new dimensions calculated through the basis functions

z = $\phi$(x) where zj = $\phi$j(x), j = 1, . . . , k

mapping from the d-dimensional x space to the k-dimensional z space where we write the discriminant as

g(z) = wT z g(x) = wT $\phi$(x)

k (13.20) = wj$\phi$j(x) j=1

where we do not use a separate w0; we assume that z1 = $\phi$1(x) $\equiv$1. Generally, k is much larger than d and k may also be larger than N, and there

lies the advantage of using the dual form whose complexity depends on N, whereas if we used the primal it would depend on k. We also use the more general case of the soft margin hyperplane here because we have no guarantee that the problem is linearly separable in this new space. The problem is the same

(13.21) Lp = 2$\parallel$w$\parallel$2 + C $\xi$t t

except that now the constraints are defined in the new space

(13.22) rtwT $\phi$(xt) $\geq$1 $-$$\xi$t

The Lagrangian is

(13.23) Lp = 2$\parallel$w$\parallel$2 + C $\xi$t $-$ $\alpha$t rtwT $\phi$(xt) $-$1 + $\xi$t $-$ $\mu$t$\xi$t t t t

When we take the derivatives with respect to the parameters and set them to 0, we get

$\partial$Lp (13.24) = w = $\alpha$trt$\phi$(xt) $\partial$w t $\partial$Lp (13.25) = C $-$$\alpha$t $-$$\mu$t = 0 $\partial$$\xi$t

The dual is now (13.26) Ld = $\alpha$t $-$1 $\alpha$t$\alpha$srtrs$\phi$(xt)T $\phi$(xs) 2 t t s

subject to

$\alpha$trt = 0 and 0 $\leq$$\alpha$t $\leq$C, $\forall$t t

The idea in kernel machines is to replace the inner product of bakernel function sis functions, $\phi$(xt)T $\phi$(xs), by a kernel function, K(xt, xs), between instances in the original input space. So instead of mapping two instances xt and xs to the z-space and doing a dot product there, we directly apply the kernel function in the original space.

(13.27) Ld = $\alpha$t $-$1 $\alpha$t$\alpha$srtrsK(xt, xs) 2 t t s

The kernel function also shows up in the discriminant

g(x) = wT $\phi$(x) = $\alpha$trt$\phi$(xt)T $\phi$(x) t (13.28) = $\alpha$trtK(xt, x) t

This implies that if we have the kernel function, we do not need to map it to the new space at all. Actually, for any valid kernel, there does exist a corresponding mapping function, but it may be much simpler to use K(xt, x) rather than calculating $\phi$(xt), $\phi$(x) and taking the dot product. kernelization Many algorithms have been kernelized, as we will see in later sections, and that is why we have the name “kernel machines.” Gram matrix The matrix of kernel values, K, where Kts = K(xt, xs), is called the Gram matrix, which should be symmetric and positive semidefinite. Recently, it has become standard practice in sharing data sets to have available only the K matrices without providing xt or $\phi$(xt). Especially in bioinformatics or natural language processing applications where x (or $\phi$(x)) has hundreds or thousands of dimensions, storing/downloading the N $\times$ N matrix is much cheaper (Vert, Tsuda, and Schölkopf 2004); this, however, implies that we can use only those available for training/testing and cannot use the trained model to make predictions outside this data set.

\section{Vectorial Kernels}

The most popular, general-purpose kernel functions are

\begin{itemize}
\item polynomials of degree q:
\end{itemize}

(13.29) K(xt, x) = (xTxt + 1)q

where q is selected by the user. For example, when q = 2 and d = 2,

K(x, y) = (xTy + 1)2 = (x1y1 + x2y2 + 1)2 = 1 + 2x1y1 + 2x2y2 + 2x1x2y1y2 + x21y21 + x22y22

corresponds to the inner product of the basis function (Cherkassky and Mulier 1998):

$\sqrt{}$ $\sqrt{}$ $\sqrt{}$ $\phi$(x) = [1, 2x1, 2x2, 2x1x2, x21, x22]T

1.5 $-$1

0.5

0 0.5 1 1.5 2

Figure 13.4 The discriminant and margins found by a polynomial kernel of degree 2. Circled instances are the support vectors.

An example is given in figure 13.4. When q = 1, we have the linear kernel that corresponds to the original formulation.

\begin{itemize}
\item radial-basis functions:
\end{itemize}

$-$x$\parallel$2 (13.30) K(xt, x) = exp $-$$\parallel$xt 2s2

defines a spherical kernel as in Parzen windows (chapter 8) where xt is the center and s, supplied by the user, defines the radius. This is also similar to radial basis functions that we discuss in chapter 12.

An example is shown in figure 13.5 where we see that larger spreads smooth the boundary; the best value is found by cross-validation. Note that when there are two parameters to be optimized using crossvalidation, for example, here C and s2, one should do a grid (factorial) search in the two dimensions; we will discuss methods for searching the best combination of such factors in section 19.2.

One can have a Mahalanobis kernel, generalizing from the Euclidean distance: " \# (13.31) K(xt, x) = exp $-$1 $-$x)T S$-$1(xt $-$x) 2(xt

(a) s2=2 (b) s2=0.5 2 2

$-$1

1 1 1 $-$1

0 0 0 1 2 0 1 2

(c) s2=0.25 (d) s2=0.1 2 2

$-$1 1

1 1

$-$1

0 $-$1 $-$1 0 0 1 2 0 1 2

Figure 13.5 The boundary and margins found by the Gaussian kernel with different spread values, s2. We get smoother boundaries with larger spreads.

where S is a covariance matrix. Or, in the most general case,

x)(13.32) K(xt, x) = exp $-$D(xt, 2s2

for some distance function D(xt, x).

\begin{itemize}
\item sigmoidal functions:
\end{itemize}

(13.33) K(xt, x) = tanh(2xT xt + 1)

where tanh($\cdot$) has the same shape with sigmoid, except that it ranges between $-$1 and +1. This is similar to multilayer perceptrons that we discussed in chapter 11.

\section{Defining Kernels}

It is also possible to define application-specific kernels. Kernels are generally considered to be measures of similarity in the sense that K(x, y) takes a larger value as x and y are more “similar,” from the point of view of the application. This implies that any prior knowledge we have regarding the application can be provided to the learner through appropriately defined kernels---“kernel engineering”---and such use of kernels can be seen as another example of a “hint” (section 11.8.4). There are string kernels, tree kernels, graph kernels, and so on (Vert, Tsuda, and Schölkopf 2004), depending on how we represent the data and how we measure similarity in that representation. For example, given two documents, the number of words appearing in both may be a kernel. Let us say D1 and D2 are two documents and bag of words one possible representation is called bag of words where we predefine M words relevant for the application, and we define $\phi$(D1) as the Mdimensional binary vector whose dimension i is 1 if word i appears in D1 and is 0 otherwise. Then, $\phi$(D1)T $\phi$(D2) counts the number of shared words. Here, we see that if we directly define and implement K(D1, D2) as the number of shared words, we do not need to preselect M words and can use just any word in the vocabulary (of course, after discarding uninformative words like “of,” “and,” etc.) and we would not need to generate the bag-of-words representation explicitly and it would be as if we allowed M to be as large as we want. Sometimes---for example, in bioinformatics applications---we can calculate a similarity score between two objects, which may not necessarily be positive semidefinite. Given two strings (of genes), a kernel measures the edit distance edit distance, namely, how many operations (insertions, deletions, substitutions) it takes to convert one string into another; this is also called alignment alignment. In such a case, a trick is to define a set of M templates and represent an object as the M-dimensional vector of scores to all the templates. That is, if mi, i = 1, . . . , M are the templates and s(xt, mi) is the score between xt and mi, then we define

$\phi$(xt) = [s(xt, m1), s(xt, m2), . . . , s(xt, mM)]T

empirical kernel and we define the empirical kernel map as map K(xt, xs) = $\phi$(xt)T $\phi$(xs)

which is a valid kernel.

Sometimes, we have a binary score function; for example, two proteins may interact or not, and we want to be able to generalize from this to scores for two arbitrary instances. In such a case, a trick is to define a graph where the nodes are the instances and two nodes are linked if they interact, that is, if the binary score returns 1. Then we say that two nodes that are not immediately linked are “similar” if the path between them is short or if they are connected by many paths. This converts pairwise local interactions to a global similarity measure, rather like defining a geodesic diffusion kernel distance used in Isomap (section 6.7), and it is called the diffusion kernel. If p(x) is a probability density, then

K(xt, x) = p(xt)p(x)

is a valid kernel. This is used when p(x) is a generative model for x measuring how likely it is that we see x. For example, if x is a sequence, p(x) can be a hidden Markov model (chapter 15). With this kernel, K(xt, x) will take a high value if both xt and x are likely to have been generated by the same model. It is also possible to parametrize the generative model as Fisher kernel p(x|$\theta$) and learn $\theta$ from data; this is called the Fisher kernel (Jaakkola and Haussler 1998).

\section{Multiple Kernel Learning}

It is possible to construct new kernels by combining simpler kernels. If K1(x, y) and K2(x, y) are valid kernels and c a constant, then ( |( cK1(x, y) (13.34) K(x, y) = K1(x, y) + K2(x, y) |( K1(x, y) $\cdot$ K2(x, y)

are also valid. Different kernels may also be using different subsets of x. We can therefore see combining kernels as another way to fuse information from different sources where each kernel measures similarity according to its domain. When we have input from two representations A and B

KA(xA, yA) + KB(xB, yB) = $\phi$A(xA)T$\phi$A(yA) + $\phi$B(xB)T $\phi$B(yB) = $\phi$(x)T $\phi$(y)

(13.35) = K(x, y)

where x = [xA, xB] is the concatenation of the two representations. That is, taking a sum of two kernels corresponds to doing a dot product in the concatenated feature vectors. One can generalize to a number of kernels

m (13.36) K(x, y) = Ki(x, y) i=1

which, similar to taking an average of classifiers (section 17.4), this time averages over kernels and frees us from the need to choose one particular kernel. It is also possible to take a weighted sum and also learn the weights from data (Lanckriet et al. 2004; Sonnenburg et al. 2006):

m (13.37) K(x, y) = $\eta$iKi(x, y) i=1

subject to $\eta$i $\geq$0, with or without the constraint of i $\eta$i = 1, respectively multiple kernel known as convex or conic combination. This is called multiple kernel learning learning where we replace a single kernel with a weighted sum. The single kernel objective function of equation 13.27 becomes (13.38) Ld = $\alpha$t $-$1 $\alpha$t$\alpha$srtrs $\eta$iKi(xt, xs) 2 t t s i

which we solve for both the support vector machine parameters $\alpha$t and the kernel weights $\eta$i. Then, the combination of multiple kernels also appear in the discriminant

(13.39) g(x) = $\alpha$trt $\eta$iKi(xt, x) t i

After training, $\eta$i will take values depending on how the corresponding kernel Ki(xt, x) is useful in discriminating. It is also possible to localize kernels by defining kernel weights as a parameterized function of the input x, rather like the gating function in mixture of experts (section 17.8)

(13.40) g(x) = $\alpha$trt $\eta$i(x|$\theta$i)Ki(xt, x) t i

and the gating parameters $\theta$i are learned together with the support vector machine parameters (Gönen and Alpaydin 2008). When we have information coming from multiple sources in different representations or modalities---for example, in speech recognition where we may have both acoustic and visual lip image---the usual approach is to feed them separately to different classifiers and then fuse the decisions;

we will discuss methods for this in detail in chapter 17. Combining multiple kernels provides us with another way of integrating input from multiple sources, where there is a single classifier that uses different kernels for inputs of different sources, for which there are different notions of similarity (Noble 2004). The localized version can then seen be an extension of this where we can choose between sources, and hence similarity measures, depending on the input.

\section{Multiclass Kernel Machines}

When there are K > 2 classes, the straightforward, one-vs.-all way is to define K two-class problems, each one separating one class from all other classes combined and learn K support vector machines gi(x), i = 1, . . . , K. That is, in training gi(x), examples of Ci are labeled +1 and examples of Ck, k /= i are labeled as $-$1. During testing, we calculate all gi(x) and choose the maximum. Platt (1999) proposed to fit a sigmoid to the output of a single (2-class) SVM output to convert to a posterior probability. Similarly, one can train one layer of softmax outputs to minimize cross-entropy to generate K > 2 posterior probabilities (Mayoraz and Alpaydin 1999):

K (13.41) yi(x) = vijfj(x) + vi0 j=1

where fj(x) are the SVM outputs and yi are the posterior probability outputs. Weights vij are trained to minimize cross-entropy. Note, however, that as in stacking (section 17.9), the data on which we train vij should be different from the data used to train the base SVMs fj(x), to alleviate overfitting. Instead of the usual approach of building K two-class SVM classifiers to separate one from all the rest, as with any other classifier, one can build K(K $-$1)/2 pairwise classifiers (see also section 10.4), each gij(x) taking examples of Ci with the label +1, examples of Cj with the label $-$1, and not using examples of the other classes. Separating classes in pairs is normally expected to be an easier job, with the additional advantage that because we use less data, the optimizations will be faster, noting however that we have O(K2) discriminants to train instead of O(K). In the general case, both one-vs.-all and pairwise separation are special ECOC cases of the error-correcting output codes that decompose a multiclass

problem to a set of two-class problems (Dietterich and Bakiri 1995) (see also section 17.6). SVMs being two-class classifiers are ideally suited to this (Allwein, Schapire, and Singer 2000), and it is also possible to have an incremental approach where new two-class SVMs are added to better separate pairs of classes that are confused, to ameliorate a poor ECOC matrix (Mayoraz and Alpaydin 1999). Another possibility is to write a single multiclass optimization problem involving all classes (Weston and Watkins 1998):

K (13.42) min $\parallel$wi$\parallel$2 + C $\xi$ti 2 i=1 i t

subject to

wzt xt + wzt0 $\geq$wixt + wi0 + 2 $-$$\xi$ti , $\forall$i /= zt and $\xi$ti $\geq$0

where zt contains the class index of xt. The regularization terms minimizes the norms of all hyperplanes simultaneously, and the constraints are there to make sure that the margin between the actual class and any other class is at least 2. The output for the correct class should be at least +1, the output of any other class should be at least $-$1, and the slack variables are defined to make up any difference. Though this looks neat, the one-vs.-all approach is generally preferred because it solves K separate N variable problems whereas the multiclass formulation uses K $\cdot$ N variables.

\section{Kernel Machines for Regression}

Now let us see how support vector machines can be generalized for regression. We see that the same approach of defining acceptable margins, slacks, and a regularizing function that combines smoothness and error is also applicable here. We start with a linear model, and later on we see how we can use kernel functions here as well:

f(x) = wTx + w0

In regression proper, we use the square of the difference as error:

e2(rt, f(xt)) = [rt $-$f(xt)]2

$-$8 $-$6 $-$4 $-$2 0 2 4 6 8

Figure 13.6 Quadratic and $\epsilon$-sensitive error functions. We see that $\epsilon$-sensitive error function is not affected by small errors and also is affected less by large errors and thus is more robust to outliers.

whereas in support vector regression, we use the $\epsilon$-sensitive loss function:

0 if |rt $-$f(xt)| < $\epsilon$ (13.43) e$\epsilon$(rt, f(xt)) = |rt $-$f(xt)| $-$$\epsilon$ otherwise

which means that we tolerate errors up to $\epsilon$ and also that errors beyond have a linear effect and not a quadratic one. This error function is thererobust regression fore more tolerant to noise and is thus more robust (see figure 13.6). As in the hinge loss, there is a region of no error, which causes sparseness. Analogous to the soft margin hyperplane, we introduce slack variables to account for deviations out of the $\epsilon$-zone and we get (Vapnik 1995)

+ + $\xi$t$-$) (13.44) min 2$\parallel$w$\parallel$2 + C ($\xi$t t

subject to

rt $-$(wTx + w0) $\leq$ $\epsilon$ + $\xi$t+ (wT x + w0) $-$rt $\leq$ $\epsilon$ + $\xi$t$-$ $\xi$t+, $\xi$t$-$ $\geq$ 0

where we use two types of slack variables, for positive and negative deviations, to keep them positive. Actually, we can see this as two hinges

added back to back, one for positive and one for negative slacks. This formulation corresponds to the $\epsilon$-sensitive loss function given in equation 13.43. The Lagrangian is

+ + $\xi$t$-$) Lp = 2$\parallel$w$\parallel$2 + C ($\xi$t t $-$ $\alpha$t+ $\epsilon$ + $\xi$t+ $-$rt + (wTx + w0) t $-$ $\alpha$t$-$ $\epsilon$ + $\xi$t$-$+ (wTx + w0) $-$rt t (13.45) $-$ ($\mu$t+$\xi$t+ + $\mu$t$-$$\xi$t$-$) t

Taking the partial derivatives, we get

$\partial$Lp (13.46) = w $-$ ($\alpha$t+ $-$$\alpha$t$-$)xt = 0 $\Rightarrow$w = ($\alpha$t+ $-$$\alpha$t$-$)xt $\partial$w t t $\partial$Lp (13.47) = ($\alpha$t+ $-$$\alpha$t$-$)xt = 0 $\partial$w0 t $\partial$Lp (13.48) = C $-$$\alpha$t+ $-$$\mu$t+ = 0 $\partial$$\xi$t+ $\partial$Lp (13.49) = C $-$$\alpha$t$-$$-$$\mu$t$-$= 0 $\partial$$\xi$t$-$

The dual is Ld = $-$1 ($\alpha$t+ $-$$\alpha$t$-$)($\alpha$s+ $-$$\alpha$s$-$)(xt)Txs 2 t s (13.50) $-$$\epsilon$ ($\alpha$t+ + $\alpha$t$-$) $-$ rt($\alpha$t+ $-$$\alpha$t$-$) t t

subject to

0 $\leq$$\alpha$t+ $\leq$C , 0 $\leq$$\alpha$t$-$$\leq$C , ($\alpha$t+ $-$$\alpha$t$-$) = 0 t

Once we solve this, we see that all instances that fall in the tube have $\alpha$t+ = $\alpha$t$-$= 0; these are the instances that are fitted with enough precision (see figure 13.7). The support vectors satisfy either $\alpha$t+ > 0 or $\alpha$t$-$> 0 and are of two types. They may be instances that are on the boundary of the tube (either $\alpha$t+ or $\alpha$t$-$is between 0 and C), and we use these to calculate w0. For example, assuming that $\alpha$t+ > 0, we have rt = xTxt + w0 + $\epsilon$. Instances that fall outside the $\epsilon$-tube are of the second type; these are

2.4 (b)

(c)

2.2

(a)

1.8

1.6

1.4

1.2 0 2 4 6 8

Figure 13.7 The fitted regression line to data points shown as crosses and the $\epsilon$tube are shown (C = 10, $\epsilon$ = 0.25). There are three cases: In (a), the instance is in the tube; in (b), the instance is on the boundary of the tube (circled instances); in (c), it is outside the tube with a positive slack, that is, $\xi$t+ > 0 (squared instances). (b) and (c) are support vectors. In terms of the dual variable, in (a), $\alpha$t+ = 0, $\alpha$t$-$= 0, in (b), $\alpha$t+ < C, and in (c), $\alpha$t+ = C.

instances for which we do not have a good fit ($\alpha$t+ = C), as shown in figure 13.7. Using equation 13.46, we can write the fitted line as a weighted sum of the support vectors:

(13.51) f(x) = wT x + w0 = ($\alpha$t+ $-$$\alpha$t$-$)(xt)T x + w0 t

Again, the dot product (xt)Txs in equation 13.50 can be replaced with a kernel K(xt, xs), and similarly (xt)T x be replaced with K(xt, x) and we can have a nonlinear fit. Using a polynomial kernel would be similar to fitting a polynomial (figure 13.8), and using a Gaussian kernel (figure 13.9) would be similar to nonparametric smoothing models (section 8.6) except that because of the sparsity of solution, we would not need the whole training set but only a subset. There is also an equivalent $\nu$-SVM formulation for regression (Schölkopf et al. 2000), where instead of fixing $\epsilon$, we fix $\nu$ to bound the fraction of support vectors. There is still a need for C though.

3.5

2.5

1.5

0 2 4 6 8

Figure 13.8 The fitted regression line and the $\epsilon$-tube using a quadratic kernel are shown (C = 10, $\epsilon$ = 0.25). Circled instances are the support vectors on the margins, squared instances are support vectors which are outliers.

(a) s2 = 5

0 2 4 6 8

(b) s2 = 0.1 2.5

1.5

0 2 4 6 8

Figure 13.9 The fitted regression line and the $\epsilon$-tube using a Gaussian kernel with two different spreads are shown (C = 10, $\epsilon$ = 0.25). Circled instances are the support vectors on the margins, and squared instances are support vectors that are outliers.

\section{One-Class Kernel Machines}

Support vector machines, originally proposed for classification, are extended to regression by defining slack variables for deviations around the regression line, instead of the discriminant. We now see how SVM can be used for a restricted type of unsupervised learning, namely, for estimating regions of high density. We are not doing a full density estimation; rather, we want to find a boundary (so that it reads like a classification problem) that separates volumes of high density from volumes of low density (Tax and Duin 1999). Such a boundary can then be used for novoutlier detection elty or outlier detection. This is also called one-class classification. one-class We consider a sphere with center a and radius R that we want to enclassification close as much as possible of the density, measured empirically as the enclosed training set percentage. At the same time, trading offwith it, we want to find the smallest radius (see figure 13.10). We define slack variables for instances that lie outside (we only have one type of slack variable because we have examples from one class and we do not have any penalty for those inside), and we have a smoothness measure that is proportional to the radius:

(13.52) min R2 + C $\xi$t t

subject to

$\parallel$xt $-$a$\parallel$2 $\leq$R2 + $\xi$t and $\xi$t $\geq$0, $\forall$t

Adding the constraints, we get the Lagrangian, which we write keeping in mind that $\parallel$xt $-$a$\parallel$2 = (xt $-$a)T(xt $-$a): \% \& (13.53) Lp = R2 + C $\xi$t $-$ $\alpha$t R2 + $\xi$t $-$ (xt)T xt $-$2aT xt + aT a $-$ $\gamma$t$\xi$t t t t

with $\alpha$t $\geq$0 and $\gamma$t $\geq$0 being the Lagrange multipliers. Taking the derivative with respect to the parameters, we get

$\partial$L (13.54) = 2R $-$2R $\alpha$t = 0 $\Rightarrow$ $\alpha$t = 1 $\partial$R t t $\partial$L (13.55) = $\alpha$t(2xt $-$2a) = 0 $\Rightarrow$a = $\alpha$txt $\partial$a t t $\partial$L (13.56) = C $-$$\alpha$t $-$$\gamma$t = 0 $\partial$$\xi$t

1.5 (b)

(c)

(a)

0.5

0 0.5 1 1.5 2

Figure 13.10 One-class support vector machine places the smoothest boundary (here using a linear kernel, the circle with the smallest radius) that encloses as much of the instances as possible. There are three possible cases: In (a), the instance is a typical instance. In (b), the instance falls on the boundary with $\xi$t = 0; such instances define R. In (c), the instance is an outlier with $\xi$t > 0. (b) and (c) are support vectors. In terms of the dual variable, we have, in (a), $\alpha$t = 0; in (b), 0 < $\alpha$t < C; in (c), $\alpha$t = C.

Since $\gamma$t $\geq$0, we can write this last as the constraint: 0 $\leq$$\alpha$t $\leq$C. Plugging these into equation 13.53, we get the dual that we maximize with respect to $\alpha$t:

(13.57) Ld = $\alpha$t(xt)Txt $-$ $\alpha$t$\alpha$s(xt)T xs t t s

subject to

0 $\leq$$\alpha$t $\leq$C and $\alpha$t = 1 t

When we solve this, we again see that most of the instances vanish with their $\alpha$t = 0; these are the typical, highly likely instances that fall inside the sphere (figure 13.10). There are two type of support vectors with $\alpha$t > 0: There are instances that satisfy 0 < $\alpha$t < C and lie on the boundary, $\parallel$xt $-$a$\parallel$2 = R2 ($\xi$t = 0), which we use to calculate R. Instances

that satisfy $\alpha$t = C ($\xi$t > 0) lie outside the boundary and are the outliers. From equation 13.55, we see that the center a is written as a weighted sum of the support vectors. Then given a test input x, we say that it is an outlier if

$\parallel$x $-$a$\parallel$2 > R2

or

xtx $-$2aT x + aTa > R2

Using kernel functions, allow us to go beyond a sphere and define boundaries of arbitrary shapes. Replacing the dot product with a kernel function, we get (subject to the same constraints):

(13.58) Ld = $\alpha$tK(xt, xt) $-$ $\alpha$t$\alpha$sK(xt, xs) t t s

For example, using a polynomial kernel of degree 2 allows arbitrary quadratic surfaces to be used. If we use a Gaussian kernel (equation 13.30), we have a union of local spheres. We reject x as an outlier if

K(x, x) $-$2 $\alpha$tK(x, xt) + $\alpha$t$\alpha$sK(xt, xs) > R2 t t s

The third term does not depend on x and is therefore a constant (we use this as an equality to solve for R where x is an instance on the margin). In the case of a Gaussian kernel where K(x, x) = 1, the condition reduces to

$\alpha$tKG(x, xt) < Rc t

for some constant Rc, which is analogous to the kernel density estimator (section 8.2.2)---except for the sparseness of the solution---with a probability threshold Rc (see figure 13.11). There is also an alternative, equivalent $\nu$-SVM type of formulation of one-class support vector machines that uses the canonical (1/2)$\parallel$w$\parallel$2 type of smoothness (Schölkopf et al. 2001).

\section{Kernel Dimensionality Reduction}

We know from section 6.3 that principal components analysis (PCA) reduces dimensionality by projecting on the eigenvectors of the covariance

(a) s2 = 1 (a) s2 = 0.1 2 2

1.5 1.5

1 1

0.5 0.5

0 0 0 1 2 0 1 2

Figure 13.11 One-class support vector machine using a Gaussian kernel with different spreads.

matrix $\Sigma$ with the largest eigenvalues, which, if data instances are centered (E[x] = 0), can be written as XT X. In the kernelized version, we work in the space of $\phi$(x) instead of the original x and because, as usual, the dimensionality d of this new space may be much larger than the data set size N, we prefer to work with the N $\times$ N matrix XXT instead of the d $\times$ d matrix XT X. The projected data matrix is $\Phi$ = $\phi$(X), and hence we work on the eigenvectors of $\Phi$T $\Phi$ and hence of the kernel matrix K. Kernel PCA Kernel PCA uses the eigenvectors and eigenvalues of the kernel matrix and this corresponds to doing a linear dimensionality reduction in the $\phi$(x) space. When ci and $\lambda$i are the corresponding eigenvectors and eigenvalues, the projected new k-dimensional values can be calculated as \$ ztj = $\lambda$ictj, j = 1, . . . , k, t = 1, . . . , N

An example is given in figure 13.12 where we first use a quadratic kernel and then decrease dimensionality to two (out of five) using kernel PCA and implement a linear SVM there. Note that in the general case (e.g., with a Gaussian kernel), the eigenvalues do not necessarily decay and there is no guarantee that we can reduce dimensionality using kernel PCA. What we are doing here is multidimensonal scaling (section 6.5) using kernel values as the similarity values. For example, by taking k = 2, one can visualize the data in the space induced by the kernel matrix, which can give us information as to how similarity is defined by the used kernel. Linear discriminality reduction (LDA) (section 6.6) can similarly

(a) Quadratic kernel in the x space 1 $-$1

$-$1 $-$1 $-$0.5 0 0.5 1

(b) Linear kernel in the z space

$-$1 0 1

$-$2

0.9 0.95 1 1.05 1.1 1.15

Figure 13.12 Instead of using a quadratic kernel in the original space (a), we can use kernel PCA on the quadratic kernel values to map to a two-dimensional new space where we use a linear discriminant (b); these two dimensions (out of five) explain 80 percent of the variance.

be kernelized (Müller et al. 2001). In chapter 6, we discussed nonlinear dimensionality reduction methods, Isomap and LLE. In fact, by viewing the elements of the cost matrix in equation 6.47 as kernel evaluations for pairs of inputs, LLE can be seen as kernel PCA for a particular choice of kernel. The same also holds for Isomap when a kernel function is defined as a function of the geodesic distance on the graph.

\section{Notes}

The idea of generalizing linear models by mapping the data to a new space through nonlinear basis functions is old, but the novelty of support vector machines is that of integrating this into a learning algorithm whose parameters are defined in terms of a subset of data instances (the dual so-called dual representation), hence also without needing to explicitly representation

evaluate the basis functions and thereby also limiting complexity by the size of the training set; this is also true for Gaussian processes where the kernel function is called the covariance function (section 14.4). The sparsity of the solution shows the advantage over nonparametric estimators, such as k-nearest neighbor and Parzen windows, or Gaussian processes, and the flexibility to use kernel functions allows working with nonvectorial data. Because there is a unique solution to the optimization problem, we do not need any iterative optimization procedure as we do in neural networks. Because of all these reasons, support vector machines are now considered to be the best, off-the-shelf learners and are widely used in many domains, especially bioinformatics (Schölkopf, Tsuda, and Vert 2004) and natural language processing applications, where an increasing number of tricks are being developed to derive kernels (Shawe- Taylor and Cristianini 2004). The use of kernel functions implies a different data representation; we no longer define an instance (object/event) as a vector of attributes by itself, but in terms of how it is similar to or differs from other instances; this is akin to the difference between multidimensional scaling that uses a matrix of distances (without any need to know how they are calculated) and principal components analysis that uses vectors in some space. More information on support vector machines can be found in books by Vapnik (1995, 1998) and Schölkopf and Smola (2002). The chapter on SVM in Cherkassky and Mulier 1998 is a very readable introduction. Burges 1998 and Smola and Schölkopf 1998 are good tutorials on SVM classification and regression, respectively. There is a dedicated Web site http://www.kernel-machines.org and many free software packages are available, such as SVMlight (Joachims 2004) and LIBSVM (Chang and Lin 2008).

\section{Exercises}

1. Propose a filtering algorithm to find training instances that are very unlikely to be support vectors.

2. In equation 13.31, how can we estimate S?

3. In the empirical kernel map, how can we choose the templates?

4. In the localized multiple kernel of equation 13.40, propose a suitable model for $\eta$i(x|$\theta$i) and discuss how it can be trained.

5. In kernel regression, what is the relation, if any, between $\epsilon$ and noise variance?

6. In kernel regression, what is the effect of using different $\epsilon$ on bias and variance?

7. How can we use one-class SVM for classification?

8. In a setting as in figure 13.12, use kernel PCA with a Gaussian kernel.

9. Let us say we have two representations for the same object and associated with each, we have a different kernel. How can we use both to implement a joint dimensionality reduction using kernel PCA?

\section{References}

Allwein, E. L., R. E. Schapire, and Y. Singer. 2000. “Reducing Multiclass to Binary: A Unifying Approach for Margin Classifiers.” Journal of Machine Learning Research 1: 113--141.

Burges, C. J. C. 1998. “A Tutorial on Support Vector Machines for Pattern Recognition.” Data Mining and Knowledge Discovery 2: 121--167.

Chang, C.-C., and C.-J. Lin. 2008. LIBSVM: A Library for Support Vector Machines. http://www.csie.ntu.edu.tw/$\sim$cjlin/libsvm/.

Cherkassky, V., and F. Mulier. 1998. Learning from Data: Concepts, Theory, and Methods. New York: Wiley.

Cortes, C., and V. Vapnik. 1995. “Support Vector Networks.” Machine Learning 20: 273--297.

Dietterich, T. G., and G. Bakiri. 1995. “Solving Multiclass Learning Problems via Error-Correcting Output Codes.” Journal of Artificial Intelligence Research 2: 263--286.

Gönen, M., and E. Alpaydin. 2008. “Localized Multiple Kernel Learning.” In 25th International Conference on Machine Learning, ed. A. McCallum and S. Roweis, 352--359. Madison, WI: Omnipress.

Jaakkola, T., and D. Haussler. 1999. “Exploiting Generative Models in Discriminative Classifiers.” In Advances in Neural Information Processing Systems 11, ed. M. J. Kearns, S. A. Solla, and D. A. Cohn, 487-493. Cambridge, MA: MIT Press.

Joachims, T. 2004. SVMlight, http://svmlight.joachims.org.

Lanckriet, G. R. G, N. Cristianini, P. Bartlett, L. El Ghaoui, and M. I. Jordan. 2004. “Learning the Kernel Matrix with Semidefinite Programming.” Journal of Machine Learning Research 5: 27--72.

Mayoraz, E., and E. Alpaydin. 1999. “Support Vector Machines for Multiclass Classification.” In Foundations and Tools for Neural Modeling, Proceedings of IWANN’99, LNCS 1606, ed. J. Mira and J. V. Sanchez-Andres, 833--842. Berlin: Springer.

Müller, K. R., S. Mika, G. Rätsch, K. Tsuda, and B. Schölkopf. 2001. “An Introduction to Kernel-Based Learning Algorithms.” IEEE Transactions on Neural Networks 12: 181--201.

Noble, W. S. 2004. “Support Vector Machine Applications in Computational Biology.” In Kernel Methods in Computational Biology, ed. B. Schölkopf, K. Tsuda, and J.-P. Vert, 71--92. Cambridge, MA: MIT Press.

Platt, J. 1999. “Probabilities for Support Vector Machines.” In Advances in Large Margin Classifiers, ed. A. J. Smola, P. Bartlett, B. Schölkopf, and D. Schuurmans, 61--74. Cambridge, MA: MIT Press.

Schölkopf, B., J. Platt, J. Shawe-Taylor, A. J. Smola, and R. C. Williamson. 2001. “Estimating the Support of a High-Dimensional Distribution.” Neural Computation 13: 1443--1471.

Schölkopf, B., and A. J. Smola. 2002. Learning with Kernels: Support Vector Machines, Regularization, Optimization, and Beyond. Cambridge, MA: MIT Press.

Schölkopf, B., A. J. Smola, R. C. Williamson, and P. L. Bartlett. 2000. “New Support Vector Algorithms.” Neural Computation 12: 1207--1245.

Schölkopf, B., K. Tsuda, and J.-P. Vert, eds. 2004. Kernel Methods in Computational Biology. Cambridge, MA: MIT Press.

Shawe-Taylor, J., and N. Cristianini. 2004. Kernel Methods for Pattern Analysis. Cambridge, UK: Cambridge University Press.

Smola, A., and B. Schölkopf. 1998. A Tutorial on Support Vector Regression, NeuroCOLT TR-1998-030, Royal Holloway College, University of London, UK.

Sonnenburg, S., G. Rätsch, C. Schäfer, and B. Schölkopf. 2006. “Large Scale Multiple Kernel Learning.” Journal of Machine Learning Research 7: 1531-- 1565.

Tax, D. M. J., and R. P. W. Duin. 1999. “Support Vector Domain Description.” Pattern Recognition Letters 20: 1191--1199.

Vapnik, V. 1995. The Nature of Statistical Learning Theory. New York: Springer.

Vapnik, V. 1998. Statistical Learning Theory. New York: Wiley.

Vert, J.-P., K. Tsuda, and B. Schölkopf. 2004. “A Primer on Kernel Methods.” In Kernel Methods in Computational Biology, ed. B. Schölkopf, K. Tsuda, and J.-P. Vert, 35--70. Cambridge, MA: MIT Press.

Weston, J., and C. Watkins. 1998. “Multiclass Support Vector Machines.” Technical Report CSD-TR-98-04, Department of Computer Science, Royal Holloway, University of London.
